\section{Related work}
\label{sec:related-work}

\subsection{Dynamic Program slicing}
Static and dynamic program slicing are not well-studied in context of functional 
programming languages. Galois slicing has been studied before for functional languages.
In \citet{perera22} the authors use an execution trace to perform slicing of functional programs.
However their work suffers from performance problems. Further, their slicing algorithms
can be difficult to interpret. Our work builds upon theirs, improving
slicing performance, and simplifies the treatment of Galois connections. Our
method also allows for us to describe other interesting queries, such
as the linked-inputs example we describe in \secref{example}.

For a generic perspective, \citet{field98} considers dynamic dependence tracking in
term-rewriting systems. Their approach allows them to perform graphical program
slicing on a subset of Pascal, but differs from ours in that they explicitly
represent reduction steps in their dependence graphs. This is in contrast to our
work, since we do not represent this sort of information in the graph, and our
language is defined with a big-step operational semantics. \todo{This could need
rewriting}. 

\subsection{Graph Evaluation}
Whilst ours is the first work to apply graph evaluation techniques to slicing
functional programs, we are not the first to apply graph evaluation techniques
in general. The compilation of a functional program to a graph has been used
to implement incremental compilation of programs in \citet{acar02, acar06b}.
In contrast to their work, which focusses on efficiently propagating changes
for inputs to computations, we are concerned with the implementation of interactive
data exploration systems. Their work is implemented as a library for Standard ML
whilst we have implemented a language with this functionality built in. 

Term graph rewriting systems, introduced by \citet{barendregt87} have been well
studied in the implementation of functional languages. In \citet{vanEekelen97},
the graph rewriting semantics of the language Clean are studied in detail, as
opposed to using an extension of the lambda calculus. As with Acar et al,
the goal of their work is the efficient implementation of pure functional
programming languages. 